\chapter{Threats to Validity and Limitations}
\label{chap:limitations}

\section{The Rubric Is an Operational Choice}

D01--D10 is a practical framework made for this thesis. It is not a universal theory of culture. Some dimensions overlap, and another researcher could group the same issues differently.

This means that the final scores should be read as results under the Vericult rubric, not as measurements of a hidden variable called true culture.

\section{Culture Is Internally Diverse}

A country, religion, language, or ethnic group can contain many different practices and opinions. A common pattern should not be turned into a rule for every person in the group.

The verifier tries to reduce this problem by asking about scope and variation. That does not remove the risk completely. A source can still overgeneralise, and the semantic model can still interpret a tendency too strongly.

\section{Evidence Availability}

Some cultural questions are well documented online. Others are not. Legal rules and official procedures are usually easier to verify than informal customs or minority practices.

Search engines also favour some languages, institutions, and regions. English-language evidence is especially overrepresented. As a result, evidence coverage partly measures what can be found online, not only what is true in the real world.

\section{Source Quality}

No source type is always correct. An official source may be the best evidence for a legal rule but a poor source for lived minority experience. A community source may provide useful local detail but be difficult to verify independently.

Vericult therefore does not use one universal source-quality number. The evidence memo has to judge whether a source is suitable for the specific question.

\section{Evidence Memos Can Be Wrong}

The evidence memo is generated by the semantic backbone. Even when the retrieved documents are good, the model can misunderstand them, miss an important qualification, or combine sources incorrectly.

For selected cases, the final analysis should therefore read both the memo and the quoted source text. The memo itself is not ground truth.

\section{The Candidate-Blind Boundary Is Not Complete}

The retrieval stage is blind to the response only after the verification questions are generated. The question generator has already seen the selected response target.

A poorly phrased question can carry the response's assumption into search. The current design reduces direct confirmation bias, but it does not remove this indirect path.

\section{Only Three Material Targets Are Checked}

Vericult checks at most three material response targets. This keeps the process understandable and limits cost, but a long answer may contain more than three important cultural claims.

If the target extractor chooses the wrong sentences, later stages can work correctly and still miss the real problem.

\section{Dependence on One Semantic Backbone}

The final verifier uses one Qwen3.6 35B-A3B backbone for the semantic stages. This keeps the setup consistent, but it also means that context extraction, question generation, evidence synthesis, and scoring share the same model family.

The thesis therefore measures the complete Vericult system with this backbone. It does not show that the same quality would hold with every other model.

\section{REPLAY Does Not Mean Perfect Determinism}

REPLAY freezes the external evidence. It does not guarantee that every future semantic-model call will return exactly the same text. Provider software, serving settings, or low-level numerical effects can still change outputs.

The reproducibility claim is therefore limited to the frozen inputs, code, configuration, and evidence artifacts.

\section{Dependence on the Response Generator}

All final responses come from one GPT-OSS generator. A different generator could produce more refusals, shorter answers, different factual errors, or different stereotypes.

The observed outcome distribution therefore describes this frozen set of responses. It is not a general estimate of how often all LLMs are culturally inappropriate.

\section{Dataset Construction}

The three corpora were built for different purposes. PLT120 and ExternalRedteam120 contain many difficult cases by design. External120 includes prompts from existing research datasets and is more varied in task type.

For this reason, the combined percentage of inappropriate responses should not be read as a normal-use failure rate.

\section{External120 Is a Curated Subset}

External120 contains 20 fixed prompts from each of six public datasets. The repository preserves the selected rows and source IDs, but the final subset should not be described as a random sample of the original datasets.

The selection improves source diversity, but it can still introduce selection effects.

\section{ExternalRedteam120 Is a Stress Test}

ExternalRedteam120 was written to expose difficult cultural assumptions. It intentionally includes identity conflicts, minority issues, religion, historical memory, and other sensitive topics.

This makes the set useful for failure analysis, but it also means that its outcome distribution cannot represent ordinary user traffic.

\section{Language Coverage}

The final dataset contains several languages through External120, but the coverage is uneven. ExternalRedteam120 is entirely in English, and many cultural settings are still absent.

Any language-specific comparison should therefore be treated as descriptive unless the sample is large enough and balanced enough for a stronger claim.

\section{Human Reference Data}

The five-person pilot study belongs to the older Best-of-4 phase. The annotators judged different responses from the final GPT-OSS set.

Those labels cannot be reused as gold data for the 360 final responses. A new human study on the exact frozen responses would be needed for direct accuracy claims.

\section{Baseline Limits}

Skywork is a preference reward model, not a calibrated cultural classifier. A raw reward score has no universal threshold that separates culturally appropriate and inappropriate answers.

The direct LLM judge has the opposite strength and weakness: it gives a clear label and explanation, but it does not use Vericult's external evidence process.

These systems are useful comparisons, but disagreement alone does not tell us which one is correct.

\section{No Direct CARB Accuracy Claim}

CARB is closely related to the topic, but its task and data differ from the final Vericult experiment. A statement such as ``Vericult is X percentage points better than CARB'' would therefore be misleading unless both were rerun on the same task with the same reference labels.

\section{Time-Sensitive Evidence}

Some cultural and institutional information changes. Laws, policies, public attitudes, and political conditions can be different in a later year.

The experiment date and frozen evidence are therefore part of the result. A later rerun may correctly retrieve different information.

\section{Ethical Limit}

An automated cultural verifier can itself become harmful if users treat its output as cultural authority. A high score does not mean that a response represents every member of a community.

The intended use is narrower: Vericult is an evaluation and debugging tool that makes its reasoning easier to inspect. Sensitive cases can still require human or community expertise.

\section{Overall Validity Claim}

The strongest claim supported by the design is limited. Vericult can process arbitrary prompt--response text in principle, and the frozen experiment shows how it behaves on the selected 360 cases.

The thesis does not claim universal cultural coverage, perfect evidence retrieval, or model-independent accuracy. These limits are part of the interpretation of the final results.
